% Part: normal-modal-logic
% Chapter: axioms-systems
% Section: consistency

\documentclass[../../../include/open-logic-section]{subfiles}

\begin{document}

\olfileid{nml}{prf}{con}

\olsection{सुसंगतता}

!!{formula}s च्या संचांचा सुसंगतता हा महत्त्वाचा
गुणधर्म आहे. !!{formula}s च्या संचापासून
$\lfalse$ सारखा व्याघात !!{derivable} असेल,
तर तो संच विसंगत असतो; अन्यथा सुसंगत असतो.
संच विसंगत असेल, तर त्यातील सर्व !!{formula}s
एखाद्या प्रतिमानातील एका जगात एकत्र सत्य असू
शकत नाहीत. संपूर्णता प्रमेयासाठी आपण याचा
व्यत्यास सिद्ध करू: प्रत्येक सुसंगत संच
एका प्रतिमानातील एका जगात सत्य असतो; ते
प्रतिमान म्हणजे ``कॅनॉनिकल प्रतिमान.''

\begin{defn}
  संच~$\Gamma$ हा प्रणाली~$\Sigma$ च्या सापेक्ष
  \emph{सुसंगत}, किंवा आपण म्हणू त्याप्रमाणे
  $\Sigma$-सुसंगत, असेल तेव्हा आणि तेव्हाच
  $\Gamma \Proves/[\Sigma] \lfalse$.
\end{defn}

उदाहरणार्थ, $\{ \Box(p \lif q), \Box p, \lnot\Box q \}$
हा संच विधानीय तर्कशास्त्राच्या सापेक्ष सुसंगत आहे,
पण \Log{K}-सुसंगत नाही. तसेच
$\{ \Diamond p, \Box\Diamond
p \lif q, \lnot q \}$ हा संच \Log{K5}-सुसंगत नाही.

\begin{prop}\ollabel{prop:consistencyfacts}
  $\Gamma$ हा !!{formula}s चा संच असू द्या. मग:
  \begin{enumerate}
  \item $\Gamma$ हा $\Sigma$-सुसंगत असेल तेव्हा आणि
    तेव्हाच असे एखादे !!{formula}~$!A$ असते की
    $\Gamma \Proves/[\Sigma] !A$.
  \item \ollabel{prop:consistencyfacts-b}%
    $\Gamma \Proves[\Sigma] !A$ असेल तेव्हा आणि
    तेव्हाच $\Gamma \cup \{
    \lnot!A \}$ हा $\Sigma$-सुसंगत नाही.
  \item \ollabel{prop:consistencyfacts-c}%
    $\Gamma$ हा $\Sigma$-सुसंगत असेल, तर
    कोणत्याही !!{formula}~$!A$ साठी
    $\Gamma \cup \{ !A \}$ हा
    $\Sigma$-सुसंगत असतो किंवा
    $\Gamma \cup \{ \lnot!A \}$ हा
    $\Sigma$-सुसंगत असतो.
  \end{enumerate}
\end{prop}

\begin{proof}
  ही तथ्ये अभिजात विधानीय तर्कशास्त्रावरून सहज
  मिळतात. \olref{prop:consistencyfacts-c} साठी
  युक्तिवाद देऊ. प्रतिपरिवर्तनाने सिद्ध करू:
  $\Gamma \cup \{ !A \}$ आणि
  $\Gamma \cup \{ \lnot!A \}$ यांपैकी कोणताही
  संच $\Sigma$-सुसंगत नाही, असे समजा. मग
  सुसंगततेच्या व्याख्येवरून
  $\Gamma, !A \Proves[\Sigma]
  \lfalse$ आणि $\Gamma, \lnot !A \Proves[\Sigma] \lfalse$
  हे दोन्ही मिळतात. निगमन प्रमेयावरून
  $\Gamma \Proves[\Sigma] !A \to \lfalse$ आणि
  $\Gamma \Proves[\Sigma] \lnot!A \lif \lfalse$.
  पण $(!A \lif
  \lfalse) \lif ((\lnot!A \lif \lfalse) \lif \lfalse)$
  हा सर्वतः-सत्य दृष्टांत आहे; म्हणून
  \olref[prp]{prop:derivabilityfacts}\olref[prp]{prop:derivabilityfacts-ruleT}
  वरून $\Gamma \Proves[\Sigma] \lfalse$.
\end{proof}

\end{document}
